\chapter{\IfLanguageName{dutch}{Deelvragen}{Answers}}
\chapter{huidige samenwerking}
De programmeurs maken \& onderhouden het EMAV project. De onderzoekers gebruiken het om berekeningen te finetunen.
\chapter{huidige staat van documentatie}
Alvorens dieper in te gaan op de tekortkomingen van de huidige aanpak, is het noodzakelijk om de huidige staat van documentatie binnen het EMAV-web project te schetsen en deze te contrasteren met de situatie in de originele Windows Forms applicatie. Momenteel bevindt het project zich in een overgangsfase, waarbij de lessen uit het verleden de koers voor de toekomst bepalen.
\section{Lessen uit het Legacy project}
In de oorspronkelijke EMAV-applicatie was de documentatie extreem gefragmenteerd. Voor de onderzoekers bestond de documentatie voornamelijk uit lijvige Word-documenten en PDF-handleidingen die losstonden van de codebasis. Voor de programmeurs waren er pogingen ondernomen om complexe business-flows (zoals de \textit{kunstmest-flow}) te documenteren in afzonderlijke Markdown-bestanden binnen de \texttt{docs}-map van de repository.
Deze aanpak leed echter aan een chronisch gebrek aan onderhoudbaarheid. In de praktijk fungeerden deze Markdown-bestanden als een soort "schaduw-code": ze bevatten grote blokken gekopieerde C\#-logica met minimale tekstuele uitleg. Dit zorgde voor een enorme \textit{documentation drift}: zodra een formule in de rekenkern werd aangepast, werd de corresponderende uitleg in de documentatiemap vergeten. Het resultaat was een systeem waarbij de documentatie eerder voor verwarring zorgde dan voor verduidelijking, wat de afhankelijkheid van de oorspronkelijke ontwikkelaars alleen maar vergrootte.
\section{Huidige aanpak in EMAV-web: Technische README's}
In de nieuwe EMAV-web applicatie is er een bewuste verschuiving merkbaar naar "Documentation as Code". De primaire vorm van documentatie voor de ontwikkelaars bestaat momenteel uit Markdown-bestanden, zoals de centrale \texttt{README.md} in de root van de repository. Deze bestanden zijn veel directer gekoppeld aan het ontwikkelproces: ze bevatten "TODO NEXT"-lijsten, instructies voor database-migraties en architecturale diagrammen in Mermaid-formaat.
Hoewel dit een aanzienlijke verbetering is voor de IT-onboarding, blijft de kloof met de onderzoekers van ILVO gapen. De informatie in deze bestanden is nagenoeg uitsluitend technisch van aard. Voor een onderzoeker die wil begrijpen hoe een emissiefactor tot stand komt, biedt een README met instructies over \texttt{WebDbContext} migraties geen enkel inzicht. De documentatie is "ontwikkelaar-tot-ontwikkelaar" geschreven, waardoor de wetenschappelijke essentie van het project nog steeds onderbelicht blijft.
\section{XML-commentaren en de limieten van DocFX}
Binnen de nieuwe C\#-codebase wordt consequent gebruikgemaakt van \texttt{/// <summary>} XML-commentaren. Er zijn concrete plannen om \textbf{DocFX} te gebruiken om deze commentaren automatisch te publiceren als een statische website. Hoewel dit de technische transparantie verhoogt, lost het de fundamentele vraag van de onderzoeker niet op. XML-commentaren zijn gebonden aan de structuur van de klassen en methoden. Ze zijn uitstekend voor het beschrijven van input-parameters, maar ongeschikt voor het weergeven van de complexe, domeinoverschrijdende wetenschappelijke redeneringen die EMAV uniek maken.
\section{De \textquotedbl Single Source of Truth\textquotedbl \@ ontbreekt}
De huidige realiteit is dat de eigenlijke "levende" documentatie van de wetenschappelijke modellen zich nog steeds buiten de broncode bevindt, in Excel-sheets en verslagen van commissies. Er is geen enkel mechanisme dat garandeert dat de formule in de Excel-sheet van de onderzoeker exact dezelfde is als de implementatie in \texttt{ILVO.EMAV.Core}. Deze versnippering is het kernprobleem dat in dit onderzoek wordt geadresseerd. De noodzaak voor een vorm van documentatie die zowel executabel is voor de machine als leesbaar voor de wetenschapper — zoals bij literate programming — komt direct voort uit deze historische en actuele tekortkomingen.
\chapter{moeilijk te verstaande onderdelen van EMAV}
Hoewel de overstap naar een modern webframework (.NET Razor Pages) veel van de technologische beperkingen van het oude Windows Forms-project wegneemt, blijft EMAV-web een applicatie met een aanzienlijke cognitieve last voor ontwikkelaars. Deze complexiteit is niet beperkt tot de architecturale keuzes, maar is diep geworteld in de aard van het domein en de technische keuzes die nodig zijn om aan de wetenschappelijke eisen te voldoen.
\section{Domeincomplexiteit: De kloof tussen Wetenschap en Implementatie}
Het meest prominente obstakel bij het begrijpen en beheren van de codebase is niet zozeer dat de programmeurs de logica niet zouden vatten — zij zijn immers verantwoordelijk voor de vertaling van de wiskundige modellen naar werkende software — maar eerder de enorme drempel die de C\#-implementatie opwerpt voor de onderzoekers zelf. In tegenstelling tot klassieke administratieve applicaties vormt de kern van EMAV een web van onderling afhankelijke wetenschappelijke formules die gebaseerd zijn op jarenlang landbouwkundig onderzoek.
Hoewel de code voor een ontwikkelaar technisch helder en logisch gestructureerd kan zijn, fungeert ze voor een onderzoeker vaak als een "black box". De wetenschappers, die de eigenlijke domeinexperts zijn, herkennen hun eigen wiskundige modellen vaak niet onmiddellijk terug in de abstracties van een objectgeoriënteerde taal. Dit zorgt voor een fundamenteel probleem in de levenscyclus van het project: wanneer een onderzoeker een formule wil finetunen of een nieuwe parameter wil testen, is hij volledig afhankelijk van de IT-afdeling. Het maken van zelfs de kleinste aanpassingen vereist een cyclus van overleg, implementatie en validatie door de programmeur, wat experimentatie en snelle iteratie ernstig belemmert. De complexiteit zit dus niet in de onbegrijpelijkheid van de code voor IT, maar in het gebrek aan autonomie voor de wetenschapper binnen een puur technische omgeving.
\section{Hybride Databeheer en Persistentiestrategieën}
Een ander complex onderdeel is de manier waarop EMAV omgaat met data. Het systeem maakt gebruik van een hybride architectuur waarbij verschillende database-technologieën naast elkaar worden gebruikt. Metadata over gebruikers, projecten en overzichten worden opgeslagen in een centrale SQL Server-database via Entity Framework Core. De ruwe inputgegevens voor berekeningen bevinden zich echter vaak in SQLite-bestanden die per project of per onderzoeksgroep worden beheerd.
Deze scheiding is noodzakelijk om onderzoekers de nodige flexibiliteit te bieden om met eigen datasets te werken, maar het introduceert aanzienlijke complexiteit in de datatoegangslaag. Ontwikkelaars moeten navigeren tussen verschillende \textit{database providers}, rekening houden met inconsistente schema's en zorgen voor data-integriteit over de grenzen van de verschillende databases heen. Het beheer van deze persistentiestromen vereist een diepgaand inzicht in zowel de infrastructuurlaag als de domeinregels.
\section{De \textquotedbl Legacy Bridge\textquotedbl \@ en Technische Schuld}
Ondanks dat EMAV-web als een nieuw project is gestart, is het onvermijdelijk dat er logica uit de oude applicatie is overgenomen. Het hergebruiken van complexe, gevalideerde rekenmodules uit de Windows Forms-tijd is essentieel om de wetenschappelijke consistentie te garanderen. Dit creëert echter een "Legacy Bridge": een zone in de code waar moderne C\#-conventies en verouderde programmeerpatronen met elkaar in conflict komen.
Het adapteren van deze oude logica naar een moderne, asynchrone Clean Architecture vereist veel "plumbing" code. Ontwikkelaars worden geconfronteerd met methoden die niet voldoen aan de huidige standaarden voor testbaarheid of leesbaarheid, maar die vanwege hun kritieke rol in de berekeningen niet zomaar herschreven kunnen worden. Deze mix van oud en nieuw verhoogt de drempel voor nieuwe ontwikkelaars om de volledige flow van de applicatie te begrijpen.
\section{Container-Orchestratie en Omgevingsbeheer}
Tot slot vormt de integratie met Docker en Jupyter Notebooks een technische uitdaging. Zoals uitgebreid besproken in Hoofdstuk \ref{ch:setup}, start de applicatie op runtime containers op voor de onderzoekers. Dit vereist dat de webapplicatie directe controle heeft over de Docker-daemon, wat complexe beveiligingsinstellingen en netwerkconfiguraties met zich meebrengt.
De noodzaak om DLL's dynamisch in te laden in een Linux-container, poortbindingen te beheren en tokens te extraheren uit logs, voegt een laag van systeem-engineering toe aan de applicatie. Dit soort logica bevindt zich op het raakvlak tussen softwareontwikkeling en systeembeheer, wat het debugging-proces bemoeilijkt: een fout kan immers zowel in de C\#-code, de Docker-configuratie als in de Python-omgeving van de container liggen.
Om deze uiteenlopende vormen van complexiteit — variërend van de wetenschappelijke detaillering tot de technische gelaagdheid van de data — beheersbaar te houden, is een ad-hoc programmeerstijl ontoereikend. In plaats daarvan steunt de architectuur van EMAV-web op een doordacht frame van software-ontwerppatronen (\textit{design patterns}). Deze patronen dienen niet louter als theoretische exercitie, maar vormen een direct antwoord op de hierboven geschetste uitdagingen.
Zo is het \textbf{Repository patroon} een noodzakelijk instrument om de hybride datastructuur te temmen, terwijl \textbf{CQRS} en het \textbf{Mediator patroon} de cognitieve last voor ontwikkelaars verlagen door verantwoordelijkheden strikt te scheiden. Tegelijkertijd biedt het \textbf{Strategy patroon} de nodige wetenschappelijke flexibiliteit binnen de rekenkern, en fungeert \textbf{Dependency Injection} als de lijm die al deze losgekoppelde componenten op een testbare manier samenbrengt. In de volgende subsecties worden de belangrijkste patronen in detail besproken, waarbij telkens de link wordt gelegd tussen het theoretische concept en de concrete meerwaarde voor de robuustheid en onderhoudbaarheid van EMAV.
\section{Repository Pattern}
Het repository patroon fungeert als een essentiële abstractielaag tussen de domeinlogica en de diverse gegevensbronnen waarover EMAV beschikt. 
Volgens \textcite{fowlerRepository} gedraagt een repository zich als een in-memory verzameling van domeinobjecten, waarbij de technische details van de onderliggende persistentiemechanismen volledig worden geabstraheerd. 
In de context van EMAV is dit patroon cruciaal omdat het project opereert op een hybride datastructuur die over de jaren heen is geëvolueerd. 
Enerzijds is er een SQLite-database voor de opslag van ruwe inputgegevens en rekendata, en anderzijds een SQL Server-database voor de webgerelateerde metadata en gebruikersbeheer.
Binnen de codebase wordt dit gerealiseerd via een gelaagde hiërarchie van interfaces en implementaties. 
In het \texttt{ILVO.EMAV.Shared}-project bevindt zich de basisdefinitie \texttt{IBaseRepository<T>}, die algemene CRUD-operaties (\textit{Create, Read, Update, Delete}) definieert. 
Specifieke projecten breiden dit uit, zoals \texttt{IWebBaseRepository<TEntity>} in de core-laag, waarbij gebruik wordt gemaakt van generieke constraints (\texttt{where TEntity \: class, IWebBaseEntity}). 
De \texttt{ILVO.EMAV.Core}-laag communiceert uitsluitend met deze interfaces. 
Hierdoor blijft de kern van de applicatie volledig onwetend over de specifieke implementatie details, of die nu gebruikmaakt van Entity Framework Core in \texttt{ILVO.EMAV.Infrastructure.Web} of gespecialiseerde logic in de input-laag.
Deze abstractie verhoogt de onderhoudbaarheid aanzienlijk\: indien een database-schema wijzigt of indien men besluit over te stappen naar een andere opslagtechnologie, hoeven de wijzigingen enkel in de infrastructuurlaag te worden doorgevoerd. 
De businesslogica blijft onaangetast. 
Bovendien is dit patroon de hoeksteen voor kwalitatieve unit testing. 
Door repositories te \textit{mocken}, kan de businesslogica worden getoetst zonder afhankelijkheid van een actieve databaseverbinding. 
Dit wijst op een architecturale keuze die gericht is op stabiliteit op lange termijn en eenvoudige verifieerbaarheid van de code \autocite{millett2015patterns}.
\section{Command Query Responsibility Segregation (CQRS)}
CQRS is een geavanceerd architectuurpatroon waarbij de verantwoordelijkheid voor het opvragen van gegevens (Queries) strikt wordt gescheiden van de verantwoordelijkheid voor het wijzigen van de systeemtoestand (Commands). 
Volgens \textcite{fowlerCQRS} biedt deze scheiding de mogelijkheid om elk pad onafhankelijk te structureren, te optimaliseren en te schalen. 
In een complex systeem zoals EMAV, waar berekeningen uren kunnen duren en metadata-opvragingen milliseconden moeten duren, is deze scheiding bittere noodzaak.
Binnen EMAV wordt CQRS geïmplementeerd met de \texttt{LiteBus}-bibliotheek \autocite{litebus}. 
Dit resulteert in een zogenaamde "Vertical Slice Architecture" binnen de core-projecten \autocite{jimmyBogardVerticalSlice}. 
Een actie zoals het verwijderen van een berekening wordt niet afgehandeld door een generieke "CalculationService", maar door een specifieke \texttt{DeleteCalculationCommand} en de bijbehorende \texttt{DeleteCalculationCommandHandler}. 
Dit commando bevat enkel de minimale data (zoals een \texttt{Guid Id}) en de handler voert de noodzakelijke stappen uit \: het ophalen van de metadata via een repository, het verwijderen van fysieke bestanden via een \texttt{IFileManager} en het finaal verwijderen van de database-records.
Aan de query-zijde worden specifieke modellen (DTO's) geretourneerd die geoptimaliseerd zijn voor de weergave in de UI, in plaats van de volledige domeinentiteiten te "lekken". 
Deze aanpak voorkomt het ontstaan van "Gouden Klassen" die alle logica van een bepaald domein bevatten. 
In plaats daarvan is de logica gedistribieerd over honderden kleine, gefocuste klassen die elk één specifieke taak hebben. 
Dit verlaagt de cognitieve last voor ontwikkelaars die aan het project werken aanzienlijk.
\section{Mediator Pattern}
Onlosmakelijk verbonden met de CQRS-aanpak is het mediator patroon. 
Dit patroon fungeert als een centrale "verkeersleider" die de communicatie tussen verschillende objecten coördineert, waardoor directe koppelingen tussen componenten worden vermeden \autocite{gof1994}. 
In de moderne architectuur van EMAV is de \texttt{IMediator} het enige aanspreekpunt voor de presentatielaag (\texttt{ILVO.EMAV.Web}).
Een Razor Page in de UI-laag hoeft geen enkel inzicht te hebben in welke services, repositories of validators nodig zijn om een bepaalde actie uit te voeren. 
De pagina stuurt simpelweg een commando of query naar de mediator. 
De mediator zorgt er vervolgens voor dat dit bericht door de juiste "pipelines" passeert. 
Zo worden commando's automatisch eerst gevalideerd door \texttt{IValidator}-implementaties voordat ze de handler bereiken.
Dit zorgt voor een extreme vorm van ontkoppeling \: de presentatielaag heeft geen enkele directe referentie naar de infrastructuur- of datalaag. 
Dit verhoogt de flexibiliteit van het systeem enorm; handlers kunnen worden herschreven, verplaatst naar andere microservices of gedecoreerd met extra logging zonder dat er ook maar één regel code in de UI-laag hoeft te wijzigen. 
Dit patroon is een directe indicatie van de architecturale volwassenheid van het project en draagt bij aan de robuustheid tegen technologische veranderingen.
\section{Strategy Pattern}
Het strategie patroon is een gedragspatroon dat toelaat om een algoritme te selecteren op runtime en dit algoritme in te kapselen in een uitwisselbaar object \autocite{gof1994}. 
In de rekenkern van EMAV is dit patroon van fundamenteel belang om de wetenschappelijke variabiliteit van de emissiemodellen te ondersteunen. 
De kerninterface \texttt{ICalculator<Tin, Tout>} vormt de basis voor een uitgebreide familie van rekenstrategieën.
Deze architectuur is essentieel omdat EMAV berekeningen uitvoert voor diverse domeinen, zoals enterische emissies, stalemissies en mestverwerking. 
Elke berekening heeft haar eigen set formules, inputparameters en validatieregels. 
Door elke berekening als een afzonderlijke strategie (bijvoorbeeld \texttt{BerekenEmissieLandbouwbodem}) te implementeren, blijft de overkoepelende orchestratielogica clean en overzichtelijk. 
De orchestrator roept simpelweg de \texttt{DoCalc}-methode aan op de geïnjecteerde strategie.
Dit bevordert het \textit{Open-Closed Principle} (onderdeel van SOLID)\: de rekenkern staat open voor uitbreiding (men kan eenvoudig een nieuwe rekenmodule toevoegen door de interface te implementeren), maar is gesloten voor aanpassingen (de bestaande orchestratie hoeft niet gewijzigd te worden). 
Voor een project waar wetenschappers regelmatig nieuwe formules aanleveren, is dit een cruciale eigenschap die voorkomt dat het project na verloop van tijd "vastloopt" in haar eigen complexiteit.
\section{Factory Pattern}
Het factory patroon wordt aangewend om de creatie van objecten te abstraheren van de client-code. 
Dit is met name nuttig wanneer het exacte type van het aan te maken object afhankelijk is van configuratie of gebruikersinput op runtime. 
In EMAV wordt dit patroon prominent toegepast in de \texttt{IExporterFactory}.
Het systeem ondersteunt een breed scala aan exportmogelijkheden, waaronder CSV voor verdere analyse in statistische software, Excel voor rapportage en JSON voor integratie met andere systemen. 
De \texttt{ExporterFactory} neemt een \texttt{ExportType} als parameter en retourneert de juiste implementatie van de \texttt{IExporter}-interface. 
Hierdoor hoeft de aanroepende code niet te weten hoe een specifieke exporter geconfigureerd moet worden.
Deze centralisatie van creatielogica voorkomt codeduplicatie en maakt het systeem zeer modulair. 
Indien er in de toekomst ondersteuning voor een nieuw formaat moet worden toegevoegd, hoeft enkel een nieuwe klasse te worden geschreven en geregistreerd in de factory. 
De rest van de applicatie blijft volledig ongewijzigd. 
Dit illustreert de visie van het project op uitbreidbaarheid en het vermijden van "harde" koppelingen tussen componenten \autocite{gof1994}.
\section{Facade Pattern}
Gezien de enorme interne complexiteit van de EMAV-rekenkern, waarbij tientallen subsystemen gelijktijdig actief zijn, is het facade patroon onontbeerlijk. 
Een facade biedt een vereenvoudigde interface naar een complex subsysteem, waardoor de interne complexiteit wordt afgeschermd voor de gebruiker \autocite{gof1994}. 
In EMAV fungeren de belangrijkste controllers en orchestratie-services als facades.
Wanneer een gebruiker op "Start Berekening" klikt, lijkt dit voor de UI een enkele actie. 
Achter de facade vindt echter een complexe dans plaats\: inputbestanden worden gevalideerd, de juiste rekenfactoren worden uit de database opgehaald, cache-mechanismen worden geraadpleegd, meerdere berekeningsstrategieën worden in de juiste volgorde aangeroepen en de resultaten worden finaal geaggregeerd en weggeschreven. 
De facade verbergt deze interacties en biedt een stabiele API.
Dit is een erkenning van de inherente complexiteit van het domein\: in plaats van te proberen de complexiteit volledig weg te nemen, wordt ze beheersbaar gemaakt door ze te groeperen achter logische, gebruiksvriendelijke interfaces. Dit verbetert de onderhoudbaarheid omdat de interne werking van de rekenkern kan worden geoptimaliseerd zonder dat de UI-code hiervan hinder ondervindt.
\section{Dependency Injection (DI) en Inversion of Control (IoC)}
Dependency Injection is de architecturale "ruggengraat" die alle voorgaande patronen in EMAV faciliteert en verbindt. 
In plaats van dat klassen hun eigen afhankelijkheden instantiëren, worden deze afhankelijkheden van buitenaf aangeboden. Volgens \textcite{seemann2019di} is DI de enige manier om een echt modulaire en testbare software-architectuur te bouwen.
EMAV maakt optimaal gebruik van de ingebouwde DI-container van .NET. 
Elk deelproject (Core, Infrastructure, Web) bevat een eigen \texttt{DependencyInjection.cs}-bestand waarin de interne services, handlers en repositories worden geregistreerd. 
Dit bevordert een hoge mate van inkapseling\: de UI-laag hoeft enkel de DI-methode van de Core-laag aan te roepen om alle benodigde logica beschikbaar te maken.
Dit mechanisme is cruciaal voor de flexibiliteit van het project. 
Het laat toe om op een transparante manier implementaties te wisselen op basis van de omgeving of voor testdoeleinden. 
DI is de motor achter de andere patronen; zonder een goede DI-structuur zouden patronen zoals Strategy en Repository leiden tot een onhoudbare hoeveelheid "handmatige" object-creatie doorheen de applicatie.
\section{Adapter Pattern}
Het adapter patroon maakt het mogelijk dat klassen met incompatibele interfaces toch naadloos kunnen samenwerken \autocite{gof1994}. 
In de context van EMAV is dit patroon een cruciaal instrument om de kloof tussen de moderne "Clean Architecture" en de waardevolle legacy-code te overbruggen. 
Zoals eerder beschreven, bevat het project logica die overgenomen is uit een oudere Windows Forms-omgeving.
In plaats van deze oude code integraal te herschrijven, worden adapters ingezet. 
Deze adapters presenteren de oude functionaliteit aan de nieuwe applicatie via een interface die voldoet aan de nieuwe standaarden (bijvoorbeeld door gebruik te maken van \texttt{Task<T>} voor asynchroon werken). 
Hierdoor kan de nieuwe kernlogica modern en "clean" blijven, terwijl ze toch gebruikmaakt van de bewezen nauwkeurigheid van de bestaande rekenmodules. 
Dit toont een pragmatische en architecturaal verantwoorde aanpak van software-evolutie aan.
\section{Decorator Pattern en Result Pattern}
Het decorator patroon voegt dynamisch verantwoordelijkheden toe aan een object zonder de onderliggende klasse te wijzigen \autocite{gof1994}. 
In de EMAV-codebase wordt dit principe toegepast via de mediator pipeline. 
Elke actie die door het systeem stroomt, wordt "gedecoreerd" met extra functionaliteit zoals automatische validatie via \texttt{FluentValidation}. 
Hierdoor blijven de eigenlijke business-handlers "puur" en gefocust op hun specifieke taak.
Daarnaast wordt het \textbf{Result Pattern} (via de \texttt{ErrorOr} bibliotheek) consistent toegepast. 
In plaats van uitzonderingen (\textit{exceptions}) te gebruiken voor control flow bij te verwachten fouten, retourneren methoden een object dat ofwel het resultaat, ofwel een lijst van fouten bevat. 
Dit patroon, zoals beschreven door \textcite{khorikovResultPattern}, zorgt voor een zeer voorspelbare en robuuste foutafhandeling die naadloos integreert met de Clean Architecture-filosofie en functionele programmeerprincipes.
\section{Clean Architecture}
Alle voorgaande patronen komen samen in de overkoepelende Clean Architecture-structuur van EMAV. 
Volgens \textcite{martin2017clean} is het hoofddoel van deze architectuur om de businesslogica volledig te isoleren van infrastructuur, frameworks en externe invloeden. 
In de projectstructuur van EMAV is deze filosofie strikt doorgevoerd\:
\begin{itemize}
    \item \textbf{Domain Layer (\texttt{ILVO.EMAV.Data.*})\:} Bevat de pure entiteiten en domeinobjecten.
    \item \textbf{Application Layer (\texttt{ILVO.EMAV.Core})\:} Bevat de businessregels, de mediator handlers en de definities van de interfaces.
    \item \textbf{Infrastructure Layer (\texttt{ILVO.EMAV.Infrastructure.*})\:} Bevat de technische implementaties, zoals database-connectiviteit en bestandssysteem-interacties.
    \item \textbf{Presentation Layer (\texttt{ILVO.EMAV.Web})\:} De user interface (Razor Pages) die uitsluitend communiceert met de Core-laag.
\end{itemize}
Deze gelaagdheid zorgt ervoor dat EMAV een zeer hoge mate van technologische onafhankelijkheid geniet. De Clean Architecture is de ultieme uiting van de complexiteitsbeheersing binnen dit project\: het biedt een duurzaam frame waarin alle andere design patterns optimaal tot hun recht komen om samen een kwalitatief superieur softwareproduct te vormen.
\chapter{tekortschieting van documentatie}
De complexiteit van EMAV-web, zoals besproken in de voorgaande secties, stelt hoge eisen aan de documentatie van het systeem. In de praktijk blijkt echter dat de traditionele methoden voor software-documentatie tekortschieten wanneer ze worden toegepast op een hybride omgeving waar wetenschappelijk onderzoek en software-engineering elkaar ontmoeten. Dit gebrek aan effectieve documentatie is een van de hoofdoorzaken van de kloof tussen de IT-afdeling en de onderzoekers.
\section{Het probleem van \textquotedbl Documentation Drift\textquotedbl}
Een van de meest hardnekkige problemen binnen het EMAV-project is de zogenaamde \textit{documentation drift}: het fenomeen waarbij de geschreven documentatie en de feitelijke broncode na verloop van tijd uit elkaar groeien. Omdat de emissiemodellen regelmatig worden bijgewerkt op basis van nieuwe wetenschappelijke inzichten, moet de documentatie telkens handmatig worden gesynchroniseerd. In een omgeving met hoge tijdsdruk wordt dit vaak verwaarloosd, waardoor onderzoekers beslissingen baseren op verouderde informatie.
Binnen de C\#-codebase uit zich dit vaak in commentaren die niet langer overeenkomen met de berekeningslogica in de handlers of rekenmodules. Wanneer een onderzoeker vraagt "hoe wordt parameter X berekend?", moet een programmeur vaak de code induiken omdat de externe documentatie (zoals Word-bestanden of Wiki-pagina's) niet langer betrouwbaar is. Dit proces is inefficiënt en verhoogt de kans op foutieve interpretaties van de wetenschappelijke modellen.
\section{Technische focus versus Domeinfocus}
Bestaande automatische documentatietools voor .NET, zoals DocFX of Swagger, zijn primair ontworpen voor softwareontwikkelaars. Ze excelleren in het weergeven van de technische interface (API-endpoints, klassenhiërarchieën, methodesiganturen), maar falen in het uitleggen van de \textit{intentie} en de wiskundige onderbouwing van de logica.
Voor een onderzoeker is de informatie dat een methode \texttt{CalculateEmission(double factor)} heet en een \texttt{Task<double>} retourneert, van weinig waarde. Wat de onderzoeker nodig heeft, is inzicht in welke wetenschappelijke bronnen aan de basis liggen van de gebruikte coëfficiënten en hoe de formules onderling samenhangen. De huidige documentatievormen zijn te "IT-centrisch" en bieden geen abstractieniveau dat aansluit bij de leefwereld van de wetenschapper.
\section{Gebrek aan interactiviteit en executabiliteit}
Traditionele documentatie is statisch. Een onderzoeker kan een PDF-document wel lezen, maar hij kan er niet mee "spelen". In een project als EMAV is het essentieel dat onderzoekers de impact van parameterwijzigingen direct kunnen observeren. Statische documentatie dwingt de onderzoeker in een passieve rol: hij moet de logica in de tekst begrijpen, deze mentaal vertalen naar de code, en vervolgens wachten op de IT-afdeling om een aanpassing door te voeren.
Het ontbreken van een executabel aspect in de documentatie betekent dat er geen "levende" link is tussen de uitleg en de uitvoering. Dit gebrek aan interactiviteit versterkt het black-box effect van de applicatie. De onderzoekers hebben nood aan een vorm van documentatie die niet alleen beschrijft \textit{wat} de code doet, maar die hen ook in staat stelt om de logica te valideren door middel van directe interactie, zonder dat zij daarvoor de volledige ontwikkelomgeving hoeven te beheersen.
\chapter{documentatie vormen}
Om de tekortkomingen van traditionele documentatie te overbruggen, wordt er in dit onderzoek gekeken naar alternatieve vormen die de scheiding tussen uitleg en uitvoering opheffen. De meest veelbelovende benaderingen zijn geworteld in de filosofie van \textit{Literate Programming} en de moderne praktische uitwerking daarvan in de vorm van computationele notebooks. Deze vormen transformeren documentatie van een statisch bijproduct naar een actief, uitvoerbaar en centraal instrument binnen de wetenschappelijke workflow.
\section{Het fundament: Literate Programming}
Zoals geïntroduceerd door \textcite{Knuth1984}, draait Literate Programming om de premisse dat een computerprogramma in de eerste plaats geschreven moet worden als een narratief voor menselijke lezers, waarbij de machine-uitvoerbare code een secundair derivaat is. In plaats van tekstuele uitleg als secundaire commentaar toe te voegen aan de code (\textit{code with comments}), vormt bij Literate Programming de tekst het primaire raamwerk waarin de codefragmenten zijn ingebed.
De essentie van Literate Programming ligt in de twee fundamentele operaties: \textbf{tangle} en \textbf{weave}. 
\begin{itemize}
    \item \textbf{Tangle:} Deze operatie extraheert de broncode uit het bronbestand en ordent deze in een formaat dat door de compiler of interpreter kan worden uitgevoerd. Voor EMAV betekent dit dat de wetenschappelijke formules die in een narratief document zijn beschreven, direct worden omgezet naar de C\#-logica die de emissieberekeningen aandrijft. Dit garandeert dat de "as-documented" logica identiek is aan de "as-implemented" logica.
    \item \textbf{Weave:} Deze operatie genereert de leesbare documentatie (zoals een PDF of HTML-pagina) uit hetzelfde bronbestand. Hierbij worden de codefragmenten gepresenteerd in hun context, vaak vergezeld van rijke visualisaties, wiskundige notaties in \LaTeX \@ en uitgebreide kruisverwijzingen.
\end{itemize}
Voor het EMAV-project biedt dit perspectief een structurele oplossing voor de \textit{documentation drift}. Wanneer de wetenschappelijke verantwoording en de eigenlijke berekeningslogica in hetzelfde artefact zijn ondergebracht, verdwijnt de noodzaak voor handmatige synchronisatie. Het document fungeert als een "Executable Paper", waarbij elke bewering over het model direct kan worden gestaafd door de bijbehorende code uit te voeren.
\section{Vergelijking van Computationele Notebook-platforms}
Computationele notebooks zijn de moderne, interactieve erfgenamen van Knuths visie. Ze bieden een "REPL-achtige" (\textit{Read-Eval-Print Loop}) ervaring die ideaal is voor exploratief onderzoek. Voor de implementatie binnen ILVO zijn drie belangrijke platforms overwogen, elk met hun eigen sterktes en zwaktes binnen een .NET-georiënteerde omgeving.
\subsection{Jupyter Notebooks}
Jupyter \autocite{jupyternotebook} is de \textit{de facto} standaard in de data science wereld. Het blinkt uit door zijn enorme ecosysteem en de ondersteuning voor honderden "kernels". In dit project is gekozen voor Jupyter vanwege de volwassenheid van de webinterface en de uitstekende integratie met Docker-omgevingen. 
Een specifiek voordeel voor EMAV is de mogelijkheid om via de \textbf{Python.NET} bridge de C\#-kern direct aan te spreken vanuit een Python-kernel. Dit laat onderzoekers toe om Python's krachtige visualisatie-libraries (zoals Matplotlib, Seaborn of Plotly) te gebruiken op data die is gegenereerd door de robuuste .NET-berekeningsengine. Het grootste nadeel van Jupyter is de niet-lineaire uitvoervolgorde van cellen, wat discipline vereist van de onderzoeker om de reproduceerbaarheid te waarborgen \autocite{Rule2018}.
\subsection{Quarto}
Quarto \autocite{quarto2025} is een technisch meer geavanceerde opvolger die specifiek focust op wetenschappelijke publicaties en technische rapportage. Het grote voordeel van Quarto is dat het gebruikmaakt van platte tekstbestanden (Markdown), wat het versiebeheer via Git aanzienlijk vergemakkelijkt in vergelijking met de JSON-gebaseerde \texttt{.ipynb} bestanden van Jupyter. Quarto ondersteunt ook "multi-language" documenten, waarbij verschillende talen in één flow kunnen worden gecombineerd. Hoewel Quarto een sterke kandidaat is voor de finale publicatie van wetenschappelijke rapporten binnen ILVO, is de interactieve "cell-by-cell" ervaring voor de gemiddelde onderzoeker momenteel nog minder intuïtief dan bij de klassieke Jupyter interface.
\subsection{Org-mode Babel}
Org-mode \autocite{orgmode2025} biedt wellicht de meest pure vorm van Literate Programming binnen de Emacs-omgeving. Met de Babel-extensie kunnen codeblokken in nagenoeg elke taal worden uitgevoerd en kunnen resultaten direct in het document worden teruggekoppeld. Voor de doeleinden van EMAV is Org-mode echter minder geschikt als breed verspreide oplossing. De steile leercurve van Emacs vormt een te grote barrière voor de gemiddelde wetenschapper die snel een berekening wil valideren. Het blijft echter een krachtige tool voor "power users" en ontwikkelaars die een maximale controle over hun documentatieproces wensen.
\section{De Techniek achter de Verweving: C\# Interop in Notebooks}
Een cruciaal aspect van de gekozen documentatievorm is de manier waarop de complexe .NET-codebase toegankelijk wordt gemaakt in de interactieve notebook-omgeving. In plaats van de logica te herimplementeren in een scripttaal (wat de betrouwbaarheid zou ondermijnen), wordt de eigenlijke \texttt{ILVO.EMAV.Core.dll} dynamisch ingeladen.
Dit proces verloopt via een startup-script dat gebruikmaakt van de \textbf{Common Language Runtime (CLR)}. Hierdoor kunnen onderzoekers in hun notebook direct objecten instantiëren van C\#-klassen, methoden aanroepen en zelfs LINQ-queries uitvoeren op de data-modellen. Deze aanpak garandeert dat de documentatie altijd werkt op de \textit{exacte} versie van de code die ook in de productieomgeving van EMAV-web draait. Het opheffen van de technische barrière tussen de gecompileerde C\#-wereld en de interactieve data science-wereld is de belangrijkste stap in het verkleinen van de kloof tussen IT en wetenschap.
\section{Interactieve Documentatie als Leer- en Validatieplatform}
De integratie van notebooks transformeert de documentatie van een passief naslagwerk naar een actieve leeromgeving. Naast de pure berekeningslogica kunnen notebooks worden verrijkt met interactieve elementen zoals sliders voor parameter-exploratie, dynamische grafieken en directe links naar externe wetenschappelijke bronnen. 
Dit zorgt voor een gelaagde informatievoorziening die aansluit bij verschillende behoeften:
\begin{enumerate}
    \item \textbf{De Narratieve Laag:} De wetenschappelijke tekst die de context, de wiskundige afleidingen en de verantwoording van het model biedt.
    \item \textbf{De Executabele Laag:} De levende codefragmenten die de formules direct implementeren en uitvoerbaar maken.
    \item \textbf{De Visualisatielaag:} Grafieken en tabellen die de output van de code direct inzichtelijk maken en vergelijking met historische data mogelijk maken.
\end{enumerate}
Door deze drie lagen in één document te verenigen, wordt voldaan aan de behoefte van de onderzoeker aan zowel transparantie als autonomie. Het resultaat is een documentatievorm die niet alleen beschrijft hoe het systeem werkt, maar die zelf een integraal en onmisbaar onderdeel vormt van het kwaliteitsborgingsproces van het instituut.
\chapter{veilig experimenteren}
Een van de belangrijkste doelstellingen van dit onderzoek is om onderzoekers de vrijheid te geven om te experimenteren met modellen, zonder dat dit een risico vormt voor de stabiliteit van de productieomgeving. Veilig experimenteren vereist zowel technische isolatie op procesniveau als functionele vangrails op dataniveau.
\section{Technische isolatie via Docker-sandboxing}
Zoals gedetailleerd beschreven in de technische opzet (zie Hoofdstuk \ref{ch:setup}), vormt containerisatie via Docker de infrastructuurele basis voor veiligheid. Elke onderzoeker krijgt bij het opstarten van een sessie een eigen, kortstondige (\textit{ephemeral}) Jupyter-container.
Dit sandboxing-mechanisme biedt bescherming op meerdere niveaus:
\begin{itemize}
    \item \textbf{Bestandssysteem-isolatie:} De container heeft een eigen \textit{read-only} root-bestandssysteem. Wijzigingen die een onderzoeker aanbrengt aan de systeemconfiguratie van de container gaan verloren na het afsluiten van de sessie, waardoor de omgeving altijd "schoon" blijft voor de volgende run.
    \item \textbf{Resource-limieten:} Door gebruik te maken van Docker's resource-quota (CPU en geheugen) kan worden voorkomen dat een experimenteel script de volledige server platlegt. Mocht een onderzoeker per ongeluk een geheugenlek of een oneindige lus introduceren, dan zal enkel zijn individuele container worden gepauzeerd of getermineerd.
    \item \textbf{Netwerk-isolatie:} De Jupyter-container bevindt zich in een afgeschermd netwerksegment. Hij kan enkel communiceren met de noodzakelijke databronnen en de web-API van EMAV, maar heeft geen ongecontroleerde toegang tot het bredere netwerk van ILVO.
\end{itemize}
Deze isolatie verlaagt de psychologische drempel voor onderzoekers om "wat-als"-scenario's te verkennen. De angst om het "echte" systeem te breken wordt technisch weggenomen door de sandbox-architectuur.
\section{Data-integriteit en Lokale Data-clonering}
Veilig experimenteren betekent ook dat de onderzoeker moet kunnen werken met realistische data zonder de centrale metadata-database te vervuilen. Binnen de EMAV-notebooks wordt gebruikgemaakt van een strategie van \textit{local cloning}.
Wanneer een onderzoeker een analyse start, wordt er een tijdelijke kopie gemaakt van de relevante SQLite-bestanden met inputgegevens. De onderzoeker kan parameters in deze lokale kopieën naar hartenlust aanpassen, verwijderen of toevoegen om de gevoeligheid van het model te testen. De officiële database blijft hierbij onaangetast. Door de resultaten van een experimentele run direct in de notebook te vergelijken met de resultaten van het officiële model, krijgt de onderzoeker onmiddellijk \textit{visual feedback}. Dit proces van directe validatie fungeert als een functionele vangrail: fouten in de wetenschappelijke logica worden direct zichtbaar in de grafieken, nog voordat er sprake is van een formele aanvraag voor een codewijziging.
\section{De weg van Experiment naar Productie: Gecontroleerde Promotie}
Tot slot maakt de voorgestelde omgeving een cruciaal onderscheid tussen een "vrij experiment" en een "formele model-update". Een notebook dient als de zandbak waarin een nieuwe wetenschappelijke formule wordt verfijnd en gevalideerd over meerdere iteraties.
Zodra een onderzoeker overtuigd is van de meerwaarde van een wijziging, fungeert de notebook als de \textit{technische specificatie} voor de IT-afdeling. In plaats van een vaag tekstdocument, overhandigt de onderzoeker een werkend code-exemplaar. De IT-ontwikkelaar neemt vervolgens de taak op zich om deze gevalideerde logica structureel te integreren in de C\#-rekenmodules, inclusief de bijbehorende unit tests en architecturale optimalisaties. Dit schept een heldere verdeling van verantwoordelijkheden: de onderzoeker is de eigenaar van de wetenschappelijke innovatie, terwijl de IT-afdeling de bewaker blijft van de softwarekwaliteit en stabiliteit.
\chapter{nauwere verweving tussen code \& docs}
De uiteindelijke meerwaarde van de voorgestelde aanpak ligt in de fundamentele verschuiving van documentatie als een statisch artefact naar documentatie als een integraal onderdeel van de software-levenscyclus. Door de scheiding tussen deze twee werelden op te heffen, ontstaat een ecosysteem dat zowel wetenschappelijk accuraat als technisch robuust is.
\section{Vermindering van de Communicatie-overhead en Ruisonderdrukking}
Een van de meest tastbare voordelen van de nauwere verweving is de drastische daling in de communicatie-overhead. In de traditionele workflow ging veel tijd verloren aan het verduidelijking van specificaties. Onderzoekers moesten wachten op IT om te zien hoe een formule in de praktijk presteerde, waarna IT vaak moest terugkomen bij de onderzoekers omdat de wiskundige beschrijving in het Word-document ambigu was.
In de nieuwe workflow fungeert de notebook als een "gemeenschappelijke taal" (\textit{ubiquitous language}). De programmeur ziet de C\#-interop en de achterliggende types, terwijl de onderzoeker de resultaten en de narratieve tekst ziet. Dit gedeelde referentiekader minimaliseert de kans op spraakverwarring. Vragen over de interne werking van de code kunnen nu vaak in enkele minuten door de onderzoeker zelf worden beantwoord door simpelweg een code-cel in de notebook uit te voeren.
\section{Documentatie als \textquotedbl Single Source of Truth\textquotedbl}
Door de principes van Literate Programming te omarmen, wordt documentatie niet langer een "sluitpost" aan het einde van een project, maar de drijvende kracht erachter. Elke wijziging in een wetenschappelijk model wordt gelijktijdig gedocumenteerd (in tekst) en geïmplementeerd (in code) binnen hetzelfde proces.
Dit verhoogt de kwaliteit van het intellectueel eigendom van ILVO aanzienlijk. De kennis zit niet meer gefragmenteerd in de hoofden van individuele experts of in een wirwar van ongecontroleerde scripts op lokale computers. In plaats daarvan is er één traceerbare bron waar de wetenschappelijke redenering en de technische executie samenkomen. Dit is essentieel voor de reproduceerbaarheid van het onderzoek, aangezien men jaren later nog exact kan zien welke aannames tot welke codewijzigingen hebben geleid.
\section{Conclusie: Een Strategische Transformatie}
De integratie van geavanceerde software-architectuur, container-isolatie en interactieve documentatievormen transformeert de relatie tussen IT en onderzoek fundamenteel. De IT-afdeling verschuift van een "uitvoerende partij" naar een "platform-enabler". Zij faciliteren de infrastructuur en de robuuste backend, terwijl de onderzoekers de domeinlogica voeden, verfijnen en valideren met een voorheen ongekende mate van autonomie.
Deze nauwe verweving is de enige duurzame oplossing voor de complexiteitsuitdagingen van EMAV-web. Het biedt een platform waarop technologische vernieuwing en wetenschappelijke precisie niet langer met elkaar in strijd zijn, maar elkaar versterken. Het verkleinen van de kloof tussen code en documentatie is hiermee niet alleen een technisch succes, maar een strategische overwinning die de wendbaarheid en de kwaliteit van het wetenschappelijk onderzoek bij ILVO structureel verbetert.
